\chapter{Tokenization -- From Text to Vectors}
\minitoc

\begin{chapteroverview}
  \begin{itemize}
    \item Compare BPE, Unigram, and WordPiece tokenization algorithms.
    \item Master the 2026 standard 128K vocabularies for multilingual performance.
    \item Implement multi-byte character handling and special token management.
    \item Analyze the trade-offs between vocabulary size and inference latency.
  \end{itemize}
\end{chapteroverview}

\section{Subword Tokenization}

BPE \parencite{sennrich2016neural} (dominant, GPT-series, 32K--128K vocabulary), WordPiece (BERT), SentencePiece \parencite{kudo2018sentencepiece} (language-agnostic, multilingual models).

\subsection{Vocabulary Size Trade-offs}

Larger vocabularies reduce sequence length (faster attention) but increase embedding parameter count and can impair performance on rare tokens. 50K--100K is the 2025 sweet spot; some multilingual models use 200K+.

\section{Universal Tokenizers}

Up to 20.2\% higher win rates in language adaptation. 5\% improvement on unseen languages. Enables efficient post-training language expansion. Key design criterion: fertility ratio (tokens per word) should be $\leq 2$ for all target languages.

\section{Sequence Packing}

Length-aware combinatorial optimization: 40\% fewer truncations, improved modeling performance. Cross-document attention masking prevents spurious inter-document dependencies within packed batches.
